% ARGUES: every departure from the pre-registration, declared, with dates and reasons.
% DRAWS FROM: PRE_REGISTRATION_OSF.md §11 amendments 1 to 13, plus the decision docs.
% GATED ON: the 15-model a-priori power result for the #13 subsection (one slot).
%
% This appendix is the paper's main credibility asset. Several amendments correct
% the author's own preferred reading, and those are the ones worth showing in full.

\section{Deviations from the pre-registration}
\label{app:deviations}

Table~\ref{tab:amendments} lists every amendment to the pre-registration in the
order it was recorded. The final column is the one a skeptical reader should
read first: it states what evidence existed at the moment each amendment was
written, because an amendment recorded before the relevant numbers exist and one
recorded after are not the same kind of object. The wording of that column is
each amendment's own stated trigger, not a later characterization of it.

Three amendments touch the interpretation of a registered result rather than an
instrument or an interface, and those are the ones worth scrutiny: \#10, \#12
and \#13. All three state their epistemic status in the pre-registration
itself. The first two were recorded before the one-shot confirmatory fit; the
third was recorded after it, discloses a scope reduction the fit inherited, and
is the one whose timing a reader should weigh most carefully. All three are
discussed in full below.

\begin{table}[htbp]
\centering
\footnotesize
\caption{Amendments to the pre-registration (§11), in the order recorded.}
\label{tab:amendments}
\begin{tabular}{@{}r l p{0.44\textwidth} p{0.26\textwidth}@{}}
\toprule
\# & Date & What changed & Evidence in hand when recorded \\
\midrule
1 & 2026-05-25 & Legal known-real pre-2023 cutoff enforced & Round-3 design panel; no generations collected \\
2 & 2026-06-04 & Anthropic snapshot pinned under alias-only model ids & Phase B pre-spend audit \\
3 & 2026-06-04 & Verification matcher and extractor abstention guard tightened & Phase B Step 0 dress rehearsal \\
4 & 2026-06-04 & Verifier recall and cross-domain scope fixed & Step 0 hand-check evidence pass \\
5 & 2026-06-06 & Phase H LLM-judge brought forward into the Phase B pipeline & Step 0 §8 A.5 hand-check gate had failed \\
6 & 2026-06-06 & Phase B cut to Haiku only; extractor swapped; judge restricted to fire on non-REAL verifier labels; subagent dispatch carve-out & Post-gate-clear cost reconciliation \\
7 & 2026-06-07 & Batch API \texttt{custom\_id} compact-code encoding & Pre-submit rejection by the API \\
8 & 2026-06-07 & Legal citations: \texttt{title} made optional at extraction load & Extraction-load failure on 34 legal citations \\
9 & 2026-06-08 & Judge dispatch moved off \texttt{claude -p} to Max-OAuth workflow dispatch & Discovery that the subprocess billed the paid API \\
10 & 2026-06-15 & Pre-committed reading of the \texttt{citations\_expected} vs \texttt{length\_matched} contrast extended to the negative-delta case, with the contrast operationally pinned & Exploratory Phase B result in a different model population; \textbf{no Phase C number on this contrast existed} \\
11 & 2026-06-15 & Phase C open-weights checkpoints added to the \texttt{BatchKey} code table & Mechanical extension of \#7; no outcome data \\
12 & 2026-08-11 & Verifier-REAL trust step recorded as failing on Phase C; corrective re-judge of at-risk citations & Blind hand-check of the Phase C corpus; \textbf{recorded before the confirmatory fit} \\
13 & 2026-09-06 & Open-weights roster disclosed as 12 of the 23 registered checkpoints; a-priori power re-measured on the roster that ran; three registered checkpoints collected under a pre-committed rule; grouping rule for the shared Llama-3.1-70B base fixed & E1 confirmatory null already in hand; \textbf{collection rule timestamped before any new generation, grouping rule before any new checkpoint was judged} \\
14 & 2026-10-01 & Registered parametric-bootstrap percentile CIs not computed; Wald CIs reported, profile-CI escalation kept on the Wald bound & Measured bootstrap cost on the 12-model frame; \textbf{filed before the 15-model fit} \\
\bottomrule
\end{tabular}
\end{table}

\subsection{Amendment \#10: the directional reversal, and the rescue that was rejected}

Section 5 of the pre-registration committed an interpretation for the
\texttt{citations\_expected} versus \texttt{length\_matched} contrast, but it
enumerated only two of the three possible signs: $\Delta \geq +5$pp and
$|\Delta| < 5$pp. The negative case was simply unspecified. That omission became
urgent when the exploratory Phase B analysis, a single closed-API model and not
the registered confirmatory fit, came back with the ordering \emph{reversed}:
\texttt{length\_matched} at an odds ratio of 6.25 $[4.25, 9.19]$ against neutral
and \texttt{citations\_expected} at only 2.79 $[1.93, 4.04]$. On that evidence
the most likely Phase C outcome was the one case the pre-registration had no
rule for.

\paragraph{The first draft was unfalsifiable, and two review passes caught it.}
This is the part worth showing. The initial draft of the amendment supplied a
reading under which the thesis survived in all three bins: a positive delta
confirmed the citation-specific cue, a null delta left it open pending more
data, and a negative delta was to be reported as evidence of a
\emph{broader} reward-hacking mechanism. Written that way, no possible value of
$\Delta$ could have counted against the claim, which is the definition of an
unfalsifiable rule and precisely what pre-registration exists to prevent. Two
adversarial review passes over the draft rejected it. The locked version makes
$\Delta \leq -5$pp an explicit \textbf{loss} for T2's citation-specific reading:
the verbosity alternative the control was designed to test is reported as not
ruled out and possibly dominant, and whether a demand for a longer answer
constitutes reward hacking at the level the thesis chain requires is reported as
a competing hypothesis this design cannot adjudicate. The $|\Delta| < 5$pp bin
was likewise rewritten so that ``unresolved'' denotes ``not supported'' rather
than ``still supported pending more data''.

\paragraph{Why the contrast is pinned to a specific computation.}
The amendment also had to fix \emph{how} $\Delta$ is computed, because the rule
is bin-ambiguous without it. On the Phase B data the model-based
\texttt{emmeans} marginal probability contrast is about $-2.0$pp, which falls in
the $|\Delta| < 5$pp bin, while the raw unconditional rate gap on the same data
is about $-7.0$pp, which falls in the $\Delta \leq -5$pp bin. The same
observations land in different bins under different definitions of the same
symbol, so the canonical one was fixed in advance: the \texttt{emmeans} marginal
probability-scale contrast with proportional weights over cells and domains,
evaluated at the \texttt{is\_post\_trained = 1} margin of the joint E2 model
rather than on a subset refit. The restriction to post-trained checkpoints is
itself pre-specified, because base checkpoints do not follow the prompt
conditions at all and the question is only defined where the cue is read.

\paragraph{The rule did not fire.}
Phase C reversed the Phase B prior decisively and under both pipelines:
$\Delta = +28.35$pp corrected and $+12.64$pp registered, both far inside the
$\Delta \geq +5$pp bin and neither close enough to a boundary to engage the
ambiguity rule. The negative bin, which is the part of this amendment that took
the most work to get right, was never needed. We report it as a rule that was
fixed in advance and did not fire, which is not the same thing as a rule that
was never needed: it was written when the available evidence pointed the other
way.

\subsection{Amendment \#12: the verifier-REAL trust step}

This amendment was not chosen; it was triggered by the pre-registration's own
rule. Amendment \#6(c) restricted the judge to fire only on verifier-FABRICATED
and verifier-UNVERIFIABLE citations, accepting verifier-REAL at face value, and
committed that if the disagreement rate on verifier-REAL citations exceeded
8\%, a sub-amendment would be recorded. That trust step was validated at
production scale on Phase B, where the registered sub-study put the
disagreement rate at 4.3\% with a design-based 95\% CI of $[1.7, 7.0]$,
clearing the gate. On the Phase C corpus a blind hand-check found 8 of 20
verifier-REAL citations were fabricated: 40.0\%, Wilson 95\% CI
$[21.9, 61.3]$. Even the lower bound is roughly three times the registered
gate, so the rule fires.

\paragraph{Why the failure was corrected rather than only disclosed.}
The registered outcome is generation-level, so a missed fabrication changes a
result only when it is the sole bad citation in its generation. Of 25,530
citing generations, 12,920 were already scored fabricated and could not flip,
leaving 5,251 at risk. What made disclosure alone inadequate is not the level of
the error but its \emph{differential} exposure across the conditions being
contrasted. The at-risk share and mean trusted-citation load run
\texttt{citations\_expected} 26.4\% and 1.64, \texttt{abstention\_permitted}
23.5\% and 1.18, \texttt{neutral} 18.9\% and 1.10, \texttt{length\_matched}
14.6\% and 1.09. Amendment \#10 pins its bin boundaries at $\pm 5$pp on exactly
the \texttt{citations\_expected} minus \texttt{length\_matched} difference, and
a differential bias of that size could manufacture a bin verdict by itself.
Exposure by stage is by contrast nearly flat, 19.7\% to 22.7\%, so E1 is not
materially implicated. Leaving this to a disclosure would have rested the
paper's most consequential pre-committed directional call on an artifact.

\paragraph{The corrective run.}
Only verifier-REAL citations inside at-risk generations were re-judged: 6,852 of
8,860, or 77.3\%, sitting in 5,251 generations. Prompts came from the registered
renderer, so the instrument is byte-identical to the deployed judge and only the
population differs. The count is 6,852 and not the 6,784 an earlier draft
carried; that planning figure was derived at 28,266 of 31,980 judgments and was
restated upward at the completed judge, along with at-risk generations
5,194 to 5,251. All 6,852 were re-judged with no duplicates and no error rows.
The verdicts on citations the structured verifier had called REAL were REAL
4,828 (70.5\%), FABRICATED 1,836 (26.8\%) and UNVERIFIABLE 188 (2.7\%), so
29.5\% of verifier-REAL labels did not survive. That 29.5\% is a net figure over
two opposing corrections, demotions where a wrong identifier resolves to a
different real work against rescues where a wrong venue or a dead DOI still
leaves one real work uniquely identified, and it is reported as a
misclassification matrix rather than as a scalar. The 188 UNVERIFIABLE
re-judgments leave the verifier's REAL label standing, matching the main combine
rule, so 6,664 of the 6,852 are applied and no generation moves into or out of
the excluded set. At the generation level the correction moves 1,458 of 25,530
citing generations, 5.71\%, every one from clean to fabricated, taking the
absolute fabrication rate from 50.61\% to 56.32\%.

\paragraph{What it did to the registered conclusion, and what it did not.}
Both pipelines are reported side by side throughout: the registered \#6(c)
pipeline with the correction withheld, and the corrected pipeline. They agree on
the amendment \#10 bin, $\Delta = +28.35$pp corrected against $+12.64$pp
registered, both in the same bin with neither interval near a boundary. The
withheld-correction pipeline is the conservative one, so the finding does not
depend on the correction having been applied, and the trust-step failure does
not reach the registered conclusion. The residual limitation is stated rather
than argued away: this replaces blind verifier trust with the judge's own
verdict, and the same hand-check puts the judge at 80.0\% precision
$[68.2, 88.2]$ on FABRICATED. The corrected outcome therefore carries roughly
20\% residual error on that stratum rather than none. It is a substantial
improvement on 40\% blind trust, not an elimination of error. Why Phase C
differs from Phase B is not established. The two studies measure the same
quantity on different corpora, and open-weights models may produce a different
mix of failure modes, including real identifiers attached to different real
works; that is a hypothesis, not a finding.

\subsection{Amendment \#13: the roster that ran was smaller than the roster registered}

This amendment is a disclosure first and a repair second. The pre-registration
enumerates 23 open-weights checkpoints across six families and states that the
E1 and E2 power figures it reports, 97.5\% and 95.5\% at nominal $\alpha$, are
computed on exactly that roster. Phase C collected twelve checkpoints across
three families: the Tülu-3 8B and 70B cascades and the OLMo-2-13B cascade, each
complete at four stages. The 405B anchor and the entire cross-family
base-to-Instruct breadth arm were never collected, for reasons of compute, and
the reduction was never recorded as a deviation. Amendment \#11 registered the
twelve checkpoints as join codes, which is where the number first appears, but
it is an interface amendment and did not propagate to the power section. The
reduction was found on 2026-09-04, after the confirmatory fit, by reading the
registered roster back against the registry the fit had used.

\paragraph{What the reduction does to the power figures.}
The registered figures describe a design that did not run. We therefore
re-measured a-priori power on the roster that did, reusing the registered
simulation unchanged and reading the roster from the confirmatory registry
rather than retyping it, with the interaction pinned to the same declared
effect the registered run pinned. Nothing in the result depends on the observed
estimate; observed power for a null would be circular and is not reported. On
the twelve-model roster E1 has 73.0\% $[66.3, 79.0]$ power at nominal $\alpha$
and 54.0\% at the Holm floor, against 97.5\% and 90.0\% registered; E2 has
73.0\% and 51.0\%, against 95.5\% and 85.0\%. All four fail the 80\% gate the
pre-registration set for itself. The two nominal figures coincide because both
are 146 rejections in 200 simulations; the Holm-floor figures show the runs are
distinct. The consequence falls on E1, whose registered result is a null: a
design that would miss the registered effect roughly 27\% of the time at nominal
$\alpha$ and 46\% at the Holm floor cannot support ``post-training does not
raise fabrication selectively.'' The null is uninformative, which is a weaker
statement than the pre-registration expected to make and weaker in the
direction that costs a claim. Low power does not rescue T1; the registered
conclusion stays ``not supported.'' E2 rejects decisively and is unaffected in
kind.

\paragraph{The repair, and the rule that governs it.}
A ladder of intermediate rosters, powered the same way, showed that collecting
registered-but-uncollected checkpoints moves E1 across the gate: adding
Llama-3.1-70B-Instruct alone reaches 80.0\%, adding the Qwen2.5-72B pair as
well reaches 91.0\%, and no reachable roster repairs the Holm floor. Before any
new generation existed we committed the rule in the repository: collect
exactly those three checkpoints, all of them in the registered roster; stop on
the registered 80\% power gate and not on significance, with the E1 estimate
never consulted; change no hypothesis, estimand, Holm family, specification or
coding; and report either outcome, an informative null or a supported T1,
without re-opening collection. The honest sequence is that E1 returned a null,
the reduction was then discovered, realized power failed the gate, and only then
were three registered checkpoints collected. Every step after the null is a
repair toward the registered design with the effect pinned a priori; what would
make it illegitimate is choosing the roster by looking at the E1 estimate, and
the rule forecloses that. The three were generated on 2026-09-06 through the
same code path and frozen question set as the twelve, 3,600 generations each
with no errors, and extracted with the same pinned extractor.

\paragraph{The ladder double-counted a base, and the grouping rule that fixes it.}
The registered roster lists Llama-3.1-70B twice, as the Tülu-3-70B stage-0
base and as the base of the Llama breadth pair, and the ladder mirrored that by
adding it as a second family. In the data it is one checkpoint with one set of
generations, so the ladder's 14- and 16-model rungs were simulated on a
structure the real data cannot carry without duplicating 3,600 observations. We
fixed the grouping before any of the three new checkpoints was judged: family
is the shared base model, so Llama-3.1-70B-Instruct joins the Tülu-3-70B family
as a sibling post-trained checkpoint and the Qwen pair is its own family. The
roster is fifteen checkpoints in four families with no duplicated rows, and the
nested random intercept keeps its registered meaning. The a-priori power on
that roster, not the ladder's figures, is what the gate is measured against.
For E1 it is 83.0\% $[77.1, 87.9]$ at nominal $\alpha$ and 69.0\% at the Holm
floor, from 200 simulations with the effect pinned as before and no errors or
warnings: the roster clears the registered gate by three points rather than the
eleven the double-counted ladder suggested, and the Holm floor stays unrepaired,
as the rule anticipated. Collection is closed under that rule. For E2 the
figure is 83.5\% $[77.6, 88.4]$ at nominal $\alpha$ and 66.5\% at the Holm
floor, again from 200 simulations with no errors or warnings, so both
confirmatory interactions clear the nominal gate and neither clears the Holm
floor.

\paragraph{What is reported.}
Both fits. The twelve-model confirmatory result in Section~\ref{sec:results}
stands as the registered one-shot fit on the roster that ran, read against its
realized power. The fifteen-model refit, same specification, both pipelines,
with Wald intervals as described in the next subsection, is reported beside it
and read against the fifteen-model power; Section~\ref{sec:refit15} gives the
result. In short, E1's registered interaction is null again at 83\% power, with
a slightly negative point estimate, and the headline T1 claim stays unsupported;
E2's interaction holds at Holm-adjusted $p = 3 \times 10^{-31}$. The same reading of the registered roster found that
the closed-API block was never collected either, so the secondary hypothesis
H1.2b is untested rather than unsupported, as Section~\ref{sec:results}
already states. The generalization claim narrows accordingly, to the base
models and post-training recipes actually run.

\subsection{Amendment \#14: Wald intervals in place of the bootstrap}

The pre-registration asks for parametric-bootstrap percentile intervals from
10,000 resamples on the target interaction terms, reported alongside Wald
intervals and p-values. We report the Wald intervals and not the bootstrap. Each
resample refits the full confirmatory model, and we measured what that costs on
the twelve-model frame with a chunked, resumable run: about 10 to 12 hours per
400 resamples for an E1 job and about 85 hours per 400 for an E2 job, at four
workers. Across the four registered jobs (E1 and E2, each under both pipelines
of amendment \#12) that comes to about 19,200 worker-hours, roughly 17,150 of
them in E2, before the larger fifteen-model frame. Adding cores did not help in
proportion, because repeated refits on a frame this size are limited by memory
bandwidth. When the run stopped it had banked most of the E1 resamples and very
few of the E2 ones. We do not pool those partial chunks, because an interval
whose resample count was set by when the run stopped is not the registered
estimator.

No registered decision depends on the bootstrap. The Holm-adjusted p-values,
the directional checks and the amendment \#10 bin all come from the fitted
models. The registration chose percentile intervals because the interaction
estimates are close to symmetric on the logit scale at this sample size, which
is also when Wald and percentile intervals agree. The one case where the
interval method could change a directional call, a bound within 2pp of the
10pp threshold, already triggers a profile-likelihood interval; we keep that
rule and apply it to the Wald bound. This amendment was filed before the
fifteen-model fit was run.

\subsection{Sensitivity \#1 elevated to co-primary}

The registered escalation trigger is monotonicity of the pooled obscure-real
exclusion rate across stages. Pooled, the rates are 2.25\%, 2.85\%, 2.58\% and
2.21\%, which is not monotone, so the literal trigger returns false. We ran the
sensitivity as co-primary anyway, and recorded the decision before the
confirmatory fit.

The reason is that the pooled test masks the domain that motivated the check.
Legal alone \emph{is} monotone decreasing across all four stages, 7.84\%,
7.72\%, 6.24\% and 5.55\%, at roughly thirty times the academic rate, and the
pooled $n$ drowns it out. The re-coding is also unusually cheap in inferential
terms: recoding UNVERIFIABLE as fabricated excludes no generation at all, so it
dissolves the non-ignorable-missingness question rather than arguing about it,
and the analysis frame goes from 42,845 to 43,200 rows. The costs are
asymmetric. Running a sensitivity that turns out to be unnecessary costs a
paragraph; omitting one that turns out to be necessary is not recoverable after
a one-shot confirmatory fit.

This is a conservative deviation in the strict sense: it adds an analysis and
changes no registered estimand, test, hypothesis family or GLMM specification.
Both halves of it have to be stated together. Reporting only that the
sensitivity was elevated would imply the registered trigger fired, and reporting
only that the trigger returned false would imply the escalation was unmotivated.
Neither is accurate. The full argument is in
\texttt{docs/legal\_domain\_holdout\_decision.md}.

\subsection{Legal domain retained in the primary analysis}

We expected to hold the legal domain out of the primary analysis and did not.
The decision was written before the confirmatory fit and is recorded in full in
\texttt{docs/legal\_domain\_holdout\_decision.md}; the grounds, in order of
weight, are these.

First, removing legal would delete a registered test. The E1 confirmatory family
enumerates \texttt{stage\_ordinal:cellobscure\_real:domainlegal} as member 5 of
five. A holdout would not be a robustness choice, it would be the silent removal
of a pre-registered hypothesis, and it would be made at exactly the moment the
result was about to be read.

Second, the rationale that pointed toward a holdout was obsolete. It came from
our own earlier Phase B work, where legal per-generation exclusion ran 63.9\% to
90.7\% and correlated with condition at 26.8 points. Those figures describe a
verifier-only pipeline. Phase C has a judge layer, and under the combined label
legal exclusion falls to roughly 2 to 6\% with a condition-correlation near 3
points. The residual is real and is disclosed in
Section~\ref{sec:limitations}, but it is no longer disqualifying on its own. It
also matters that the Phase B precedent was an analysis choice made at the time
and never registered, so it cannot license removing a registered family member
now.

Third, the pre-registration already contains its own remedy for informative
exclusion, and it is not a holdout. It is the sensitivity \#1 escalation
described above, which is domain-agnostic, excludes nothing, and was run as
co-primary. Using the registered remedy is preferable to inventing an
unregistered one after the fact.

% Scope decision made 2026-08-12: build E3 and E4, so the §2 single-publication
% commitment is met and no scope deviation is declared. E3 was built and run
% 2026-09-05/06 on the registered scale-check arm; the deviations that entails
% are recorded below.

\subsection{E3 run on the registered scale-check arm, from quantized weights}

The pre-registration names three probe targets for E3: Tülu-3-70B-Instruct as
the primary, Tülu-3-8B as a scale check, and Llama-3.1-70B-Instruct as a
cross-pipeline check. Only the scale check was run. H1.3 is stated on the 70B
primary, so the registered hypothesis is untested and the reported result is
the registered procedure on the registered secondary model. The decision was
made on 2026-09-05, before any E3 number existed, and the reason is compute
rather than outcome: the study's remaining budget was zero dollars and the local
GPU holds 12~GB.

Two further departures follow from the same constraint, one specification
fills a gap the registration left open, and one departure comes from the data.
First, the 8B checkpoint was loaded in 4-bit NF4 quantization, because 8B
in bf16 does not fit 12~GB; the hidden states are therefore read from a
quantized copy of the weights that produced the Phase C generations. Second, the
registered difficulty proxy, the base model's log-probability of producing any
citation for the item, was estimated from Llama-3.1-8B's existing Phase C
generations, four draws per question, rather than from a dedicated scoring
pass, which leaves it with five distinct levels; the draw count is stored beside
every value. Third, the layer was selected on the probe set's own validation
split and the full 33-layer sweep is reported, so no layer was chosen on
citation data. Fourth, the registered positive class is ``fabricated plus
misgrounded,'' and no Phase C citation carries the MISGROUNDED label, so the
class is FABRICATED alone; UNVERIFIABLE citations are excluded from both
classes.

On timing: the probe, the cross-task and within-domain estimates, the layer
sweep and the agreement rule were implemented and run on 2026-09-05. The
deconfound's implementation, the proxy construction and the orientation rule
under which the difficulty probe is scored in whichever direction is larger,
was written on 2026-09-06, after the within-domain AUROC of $0.886$ was known
and before the deconfound itself was run. The orientation rule is conservative
for the fabrication claim, and the proxy's definition is the registered one; we
record the order anyway, because a guardrail written after its target number is
known is not the same object as one written before.
